% config/layout.tex — 页面布局、几何尺寸与页眉页脚
% ===================================================

% ========== 统一尺寸参数（全部手工设置，无任何自动计算）==========

\newlength{\SideMargin}  
\setlength{\SideMargin}{1.6cm}

\newlength{\TopMargin}
\setlength{\TopMargin}{3cm}

\newlength{\BottomMargin}
\setlength{\BottomMargin}{3cm}

\newlength{\PartTitleRaise}
\setlength{\PartTitleRaise}{-3cm}

\newlength{\ChapterRightMargin}
\setlength{\ChapterRightMargin}{5.5cm}

\newlength{\MarginNoteWidth}
\setlength{\MarginNoteWidth}{3.4cm}

\newlength{\MarginNoteSep}
\setlength{\MarginNoteSep}{0.6cm}

% ========== 默认几何（对称）==========
\geometry{
  left       = \SideMargin,
  right      = \SideMargin,
  top        = \TopMargin,
  bottom     = \BottomMargin,
  headheight = 15pt,
  headsep    = 8mm,
  footskip   = 12mm,
}

\raggedbottom

% ========== 几何切换命令 ==========

\newlength{\SymLeftMargin}
\setlength{\SymLeftMargin}{3cm}
\newlength{\SymRightMargin}
\setlength{\SymRightMargin}{3cm}

% 对称布局（封面、前言、目录、part cover、附录、致谢）
% 无页眉，页码居中底部（twoside 时左右分置）
\newcommand{\SymmetricGeometry}{%
  \newgeometry{
    left       = \SymLeftMargin,
    right      = \SymRightMargin,
    top        = \TopMargin,
    bottom     = \BottomMargin,
    headheight = 15pt,
    headsep    = 8mm,
    footskip   = 12mm,
  }%
  % marginnote 会缓存 textwidth；切换几何后需要同步刷新
  \edef\marginnotetextwidth{\the\textwidth}%
  \fancyhfoffset{0pt}%
  \pagestyle{symplain}%
}

\newcommand{\ChapterGeometry}{%
  \iftwosidemode
    \newgeometry{
      inner      = \SideMargin,
      outer      = \ChapterRightMargin,
      top        = \TopMargin,
      bottom     = \BottomMargin,
      marginparwidth = \MarginNoteWidth,
      marginparsep   = \MarginNoteSep,
      headheight = 15pt,
      headsep    = 8mm,
      footskip   = 12mm,
    }%
  \else
    \newgeometry{
      left       = \SideMargin,
      right      = \ChapterRightMargin,
      top        = \TopMargin,
      bottom     = \BottomMargin,
      marginparwidth = \MarginNoteWidth,
      marginparsep   = \MarginNoteSep,
      headheight = 15pt,
      headsep    = 8mm,
      footskip   = 12mm,
    }%
  \fi
  \setlength{\marginparwidth}{\MarginNoteWidth}%
  \setlength{\marginparsep}{\MarginNoteSep}%
  % 同步 marginnote 的版心基准，避免右侧注释被旧 textwidth 推出页面
  \edef\marginnotetextwidth{\the\textwidth}%
  % 页脚保持版心宽度；仅页眉扩展到右侧边注区
  \fancyfootoffset[L]{0pt}%
  \fancyfootoffset[R]{0pt}%
  \fancyheadoffset[L]{0pt}%
  \iftwosidemode
    \fancyheadoffset[OR]{\dimexpr \ChapterRightMargin - \SideMargin\relax}%
    \fancyheadoffset[EL]{\dimexpr \ChapterRightMargin - \SideMargin\relax}%
  \else
    \fancyheadoffset[R]{\dimexpr \ChapterRightMargin - \SideMargin\relax}%
  \fi
  \pagestyle{chapterfancy}%
}

% ========== 页眉页脚 ==========

% 默认页样式：无页眉
\pagestyle{fancy}
\fancyhf{}
\renewcommand{\headrulewidth}{0pt}
\renewcommand{\footrulewidth}{0pt}

% symplain：无页眉，页码底部（oneside 居中，twoside 左右分置）
\fancypagestyle{symplain}{%
  \fancyhf{}%
  \renewcommand{\headrulewidth}{0pt}%
  \iftwosidemode
    \fancyfoot[LE,RO]{\thepage}%
  \else
    \fancyfoot[C]{\thepage}%
  \fi
}

% chapterfancy：有页眉 + 页码底部
\fancypagestyle{chapterfancy}{%
  \fancyhf{}%
  \renewcommand{\headrulewidth}{0.4pt}%
  \iftwosidemode
    \fancyhead[LE]{\nouppercase{\rightmark}}%
    \fancyhead[RO]{\today}%
    \fancyfoot[LE,RO]{\thepage}%
  \else
    \fancyhead[L]{\nouppercase{\rightmark}}%
    \fancyhead[R]{\today}%
    \fancyfoot[C]{\thepage}%
  \fi
}

% plain：章首页（\chapter 自动使用），无页眉，页码底部
\fancypagestyle{plain}{%
  \fancyhf{}%
  \renewcommand{\headrulewidth}{0pt}%
  \iftwosidemode
    \fancyfoot[LE,RO]{\thepage}%
  \else
    \fancyfoot[C]{\thepage}%
  \fi
}

% ========== 行距与段落 ==========
\setstretch{1.5}
\setlength{\parskip}{1.5ex plus 0.4ex minus 0.2ex}

% ========== sidenote 命令 ==========
% 保持旁注文字大小不变，仅缩小旁注中的数学公式。
% 档位（离散）: \textstyle / \scriptstyle / \scriptscriptstyle
% 中间值（连续缩放）：在旁注中用 \sidemath{...}，可微调 0.90~0.98
\newcommand{\SidenoteMathScale}{0.8}
\newcommand{\sidemath}[1]{\scalebox{\SidenoteMathScale}{#1}}
\newcommand{\sidenote}[1]{%
  \marginnote{%
    \raggedright
    \fontsize{10pt}{12.5pt}\selectfont
    \kaiti
    \setstretch{1.1}%
    #1%
  }%
}

% ========== Part 标题格式 ==========
\ctexset{
  chapter = {
    name = {第,章},
    number = \arabic{chapter},
    format = \Huge\bfseries\filcenter,
  },
  part = {
    pagestyle = empty,
    beforeskip = -\PartTitleRaise,
    aftername = \par\vskip 8pt\noindent\rule{\textwidth}{2pt}\par\vskip 8pt,
  }
}

% 取消 \part 后自动换页，改为输出下方水平线 + 间距，以容纳局部目录
\makeatletter
\renewcommand{\@endpart}{%
  \par\vskip 10pt
  \noindent\rule{\textwidth}{2pt}%
  \par\vskip 1cm
}
\makeatother

% ========== 图片搜索路径 ==========
\graphicspath{{_asset/figures/}}



\newcommand{\myskip}{\vspace{0.5cm}}
\newcommand{\mysec}[1]{\myskip\noxmubluebf{#1}}
